\section{Part C. BFS}
\begin{frame}{BFS의 Mission}
\gmission{source에서 edge 수가 작은 layer부터 탐색하고 unweighted shortest-path distance를 계산한다.}
\statelegend
\begin{itemize}\item WHITE: undiscovered\item GRAY: discovered and queued\item BLACK: fully processed\end{itemize}
\end{frame}
\begin{frame}{BFS Pseudocode}
\small
\begin{algorithmic}[1]
\Procedure{BFS}{$G,s$}\ForAll{$v\in V$}\State $color[v]\gets W,\ dist[v]\gets\infty,\ parent[v]\gets\NIL$\EndFor
\State $color[s]\gets G,\ dist[s]\gets0$; \Call{Enqueue}{$Q,s$}
\While{$Q\ne\emptyset$}\State $u\gets\Call{Dequeue}{Q}$
\ForAll{$v\in Adj[u]$}\If{$color[v]=W$}\State $color[v]\gets G,\ dist[v]\gets dist[u]+1$
\State $parent[v]\gets u$; \Call{Enqueue}{$Q,v$}\EndIf\EndFor
\State $color[u]\gets B$\EndWhile\EndProcedure
\end{algorithmic}
\end{frame}
\begin{frame}{왜 Enqueue할 때 발견 표시하는가?}
\begin{columns}[T]\column{.48\textwidth}\begin{alertblock}{Dequeue 때 표시}두 frontier vertex가 같은 neighbor를 각각 enqueue할 수 있다.\end{alertblock}
\column{.48\textwidth}\begin{block}{Enqueue 때 표시}첫 발견이 slot을 예약한다. 각 vertex는 queue에 최대 한 번 들어간다.\end{block}\end{columns}
\gtakeaway{GRAY는 “처리 완료”가 아니라 “이미 발견되어 queue에 있음”이다.}
\end{frame}
\begin{frame}{고정 BFS 예제}
\small adjacency order:
\[
\begin{aligned}
0&:[1,2,3]&1&:[0,4]&2&:[0,4,5]&3&:[0,6]\\
4&:[1,2,7]&5&:[2]&6&:[3,7]&7&:[4,6].
\end{aligned}
\]
\centering\begin{tikzpicture}[x=1.25cm,y=.78cm]
\node[gvertex](v0)at(0,0){0};\node[gvertex](v1)at(2,1){1};\node[gvertex](v2)at(2,0){2};\node[gvertex](v3)at(2,-1){3};
\node[gvertex](v4)at(4,1){4};\node[gvertex](v5)at(4,0){5};\node[gvertex](v6)at(4,-1){6};\node[gvertex](v7)at(6,0){7};
\draw[gedge](v0)--(v1)(v0)--(v2)(v0)--(v3)(v1)--(v4)(v2)--(v4)(v2)--(v5)(v3)--(v6)(v4)--(v7)(v6)--(v7);
\end{tikzpicture}
\end{frame}
\begin{frame}{BFS Queue Trace}
\centering
\begin{tikzpicture}
\node[callout]at(0,1.4){\only<1>{start 0}\only<2>{dequeue 0; discover 1,2,3}\only<3>{dequeue 1; discover 4}\only<4>{dequeue 2; discover 5}\only<5>{dequeue 3; discover 6}\only<6>{process 4,5,6; discover 7}\only<7>{dequeue 7; discover none}};
\node[pq]at(0,.3){Queue: \only<1>{0}\only<2>{1,2,3}\only<3>{2,3,4}\only<4>{3,4,5}\only<5>{4,5,6}\only<6>{7}\only<7>{empty}};
\node[callout]at(0,-.8){Output: \only<1>{--}\only<2>{0}\only<3>{0,1}\only<4>{0,1,2}\only<5>{0,1,2,3}\only<6>{0,1,2,3,4,5,6}\only<7>{0,1,2,3,4,5,6,7}};
\end{tikzpicture}
\begin{block}{Convention}Output에는 queue에서 제거되어 처리된 vertex만 기록한다.\end{block}
\end{frame}
\begin{frame}{BFS Layer와 Parent}
\[
\begin{array}{c|cccccccc}v&0&1&2&3&4&5&6&7\\\hline
dist&0&1&1&1&2&2&2&3\\parent&-&0&0&0&1&2&3&4
\end{array}
\]
\begin{block}{Invariant}queue에서 제거되는 vertex의 distance는 감소하지 않는다. 첫 발견보다 짧은 layer는 이미 모두 처리되었으므로 첫 \(dist\)가 최단이다.\end{block}
\end{frame}
\begin{frame}{BFS 순서는 하나가 아니다}
\small\centering\begin{tabular}{lcc}\toprule&Adj[0]&가능한 앞부분\\\midrule
order A&[1,2,3]&0,1,2,3,\ldots\\
order B&[3,2,1]&0,3,2,1,\ldots\\\bottomrule\end{tabular}
\begin{itemize}\item 달라질 수 있음: 같은 layer의 traversal order, parent tree\item 유지됨: distance layer와 shortest edge-count distance\end{itemize}
\end{frame}
\begin{frame}{Path Reconstruction}
\begin{tikzpicture}[x=1.4cm]\node[gvertex](a)at(0,0){7};\node[gvertex](b)at(2,0){4};\node[gvertex](c)at(4,0){1};\node[gvertex](d)at(6,0){0};
\draw[gactive](a)--node[above]{parent}(b);\draw[gactive](b)--(c);\draw[gactive](c)--(d);\end{tikzpicture}
\[
7\leftarrow4\leftarrow1\leftarrow0\quad\Longrightarrow\quad0\to1\to4\to7.
\]
\begin{block}{Unreachable}\(dist=\infty,\ parent=\NIL\). parent chain은 source에 도달하지 않는다.\end{block}
\end{frame}
\begin{frame}{Disconnected Graph와 BFS Forest}
\begin{itemize}\item single-source BFS는 source component만 방문한다.\item 전체 component를 찾으려면 undiscovered source마다 BFS-Component를 시작하고 component id를 기록한다.\item directed graph에서 이 반복은 SCC algorithm이 아니다.\end{itemize}
\end{frame}
\begin{frame}{BFS Complexity}
\graphcomplexity{\(\Theta(V+E)\)}{\(\Theta(V^2)\)}
\begin{itemize}\item vertex enqueue/dequeue 최대 1회\item directed edge 1회, undirected adjacency entry 2회 검사\item arrays \(\Theta(V)\), queue \(O(V)\); graph storage 별도\end{itemize}
\end{frame}
\begin{frame}{Part C Checkpoint}
\begin{enumerate}\item 왜 visited를 enqueue 시 표시하는가?\item 0에서 7까지 복원된 path는?\item adjacency order가 바뀌어도 유지되는 값은?\end{enumerate}
\begin{exampleblock}{답}중복 enqueue 방지; \(0\to1\to4\to7\); shortest edge-count distance.\end{exampleblock}
\end{frame}
